\section{The scope of a common-cover reciprocity argument}
\label{sec:common-cover-reciprocity-limits}

The remaining harmonic comparison is
\begin{equation}
 \overline{B_{g,p}(i,-1)}^{(p-\chi_{-4}(p))/2}
      =(-1)^{\widetilde\psi_p(g)}
      =(-1)^{\rho_p^I(g)/2},\qquad g\in N.               \label{eq:common-cover-reciprocity-question}
\end{equation}
Here the bar denotes reduction in the residue norm-one torus.
The companion equality proves the second equality on all of \(M\),
and hence on \(N\).
The first equality compares two different coefficient primes.
The following proposition specifies what an integral common cover and
the componentwise beta identities do not, by themselves, prove.
It does not construct a counterexample in GT.

\begin{proposition}\label{prop:common-cover-homology-countermodel}
For every odd prime \(p\), the degree-four cyclic cover with
monodromies \((1,2)\) and the degree-\(p\) cyclic cover with
monodromies \((1,1)\) have a common degree-\(4p\) cover.
Its compact genus is \(2p-1\), and it has four punctures.

There is an automorphism of its full profinite integral first
homology which commutes with the deck group, preserves the
intersection form, fixes every peripheral homology class,
and is compatible with pushforward and transfer to both smaller
covers. It induces the identity on the full integral homology
of the degree-\(p\) cover, but has nontrivial harmonic quadratic
sign at coefficient \(p\) on the degree-four cover. In particular
it is the identity at coefficient two.
\end{proposition}

\begin{proof}
The common deck group is \(Q=C_4\times C_p\), with the
monodromy images
\[
                  x=(1,1),\qquad y=(2,1),\qquad z=(1,-2).
\]
Their orders are \(4p,2p,4p\), respectively. Thus the three
original punctures have \(1,2,1\) preimages. The Euler
characteristic of the open cover is \(-4p\), so
\(2-2g-4=-4p\) gives \(g=2p-1\).

Let \(a=(1,0)\in Q\), of order four, and let \(V_p\)
be the integral first homology of the open common cover over
\(\mathbf Z_p\). Since \(p\) is odd, the following decomposition
is integral and etale:
\begin{equation}
 \mathbf Z_p[C_4]
      \simeq\mathbf Z_p\times\mathbf Z_p\times\mathbf Z_p[i].
                                                        \label{eq:common-cover-c4-algebra}
\end{equation}
The factors correspond to the trivial, order-two, and primitive
order-four characters. The involution induced by \(a\mapsto a^{-1}\)
is trivial on the first two factors and complex conjugation on
the last one. This notation includes the split algebra when
\(p\equiv1\pmod4\).

The residue norm-one group in the last factor is cyclic of order
\(p-\chi_{-4}(p)\). Choose an element with nontrivial quadratic
sign, and lift it to a norm-one unit \(v\in\mathbf Z_p[i]^\times\).
Such a lift can be taken among the prime-to-\(p\) roots of unity:
in the split case lift the two reciprocal coordinates, and in
the inert case lift a root of unity of order \(p+1\).
Let \(u\in\mathbf Z_p[C_4]^\times\) be the element with
components \((1,1,v)\) in \eqref{eq:common-cover-c4-algebra}. Then
\[
                              u^\iota u=1.
\]
Let \(u(a)\) act on \(V_p\). It is an automorphism and
commutes with the entire abelian deck group. Deck transformations
preserve the intersection form, so the adjoint of \(a\) is
\(a^{-1}\). Hence \(u^\iota u=1\) proves that \(u(a)\)
preserves this form, including its induced nondegenerate form
on compact homology.

The boundary representation is the sum of the permutation
representations of the four punctures, modulo their one sum
relation. The two punctures over \(y\) have stabilizer
\(\langle(2,1)\rangle\), so their quotient deck action is
\(C_2\); the other punctures are fixed. Consequently every
boundary class lies in the first two factors of \eqref{eq:common-cover-c4-algebra}.
The element \(u\) fixes all of them, integrally.

Under the map to the degree-four cover, \(a\) becomes its
deck generator. Thus \(u(a)\) induces multiplication by
\(v\) on its compact primitive order-four summand, giving
the prescribed nontrivial harmonic sign modulo \(p\).
Under the map to the degree-\(p\) cover, \(a\) becomes one;
the induced operator is the augmentation of \(u\), which
is one. Pushforward and transfer commute with these deck
actions. A polynomial in \(a\), with \(\mathbf Z_p\)
coefficients, therefore commutes with both maps in the
corresponding equivariant sense.

Finally choose this operator at coefficient \(p\) and the
identity at every other coefficient prime. The product is
an automorphism of full profinite integral homology. It
has all the stated properties. No denominators divisible
by \(p\) occur in \eqref{eq:common-cover-c4-algebra}, so the construction does not discard
the integral structure at the coefficient prime where the
\(C_p\)-deck action is ramified.
\end{proof}

This proposition concerns homology operators. It does not lift
the operator to an automorphism of the ambient free profinite
group, to a positive peripheral automorphism, or to a solution
of the GT pentagon. It shows that an argument using only
integral homology, intersection pairings, the indicated deck
action, and the common-cover transfer maps leaves both signs
possible in \eqref{eq:common-cover-reciprocity-question}.


The componentwise formal construction in
Proposition~\ref{prop:harmonic-formal-residue-gap} gives a separate
limitation: beta associativity, unitarity, integral dyadic duplication
and the logarithmic Anderson identities do not force the first equality
in \eqref{eq:common-cover-reciprocity-question}. The companion theorem
now identifies its other side with the divided Kummer character.
Neither construction gives a full-GT element with nonzero discrepancy.
A decision requires further nonabelian compatibility or a GT witness.
